\documentclass[acmsmall,anonymous=false,nonacm]{acmart}

%% ---------------------------------------------------------------------------
%%  Verified Optimistic Extraction from the Calculus of Constructions
%%  to a Polymorphic Blame Calculus (System Fomega + blame)
%% ---------------------------------------------------------------------------

\usepackage{mathpartir}
\usepackage{listings}
\usepackage{booktabs}
\usepackage{xcolor}

\settopmatter{printacmref=false}
\setcopyright{none}

%% --- notation macros -------------------------------------------------------
\newcommand{\dyn}{\mathord{?}}
\newcommand{\kstar}{\ensuremath{\star}}
\newcommand{\karr}{\ensuremath{\Rightarrow}}
\newcommand{\Prop}{\mathsf{Prop}}
\newcommand{\Set}{\mathsf{Set}}
\newcommand{\Kind}{\mathsf{Kind}}
\newcommand{\Nat}{\mathsf{Nat}}
\newcommand{\srt}{s}
\newcommand{\cast}[4]{\langle #2 \Rightarrow #3 \rangle^{#4}\, #1}
\newcommand{\gnd}[2]{\mathsf{gnd}(#1, #2)}
\newcommand{\isgnd}[2]{\mathsf{is\_gnd}(#1, #2)}
\newcommand{\blame}[1]{\mathsf{blame}\,#1}
\newcommand{\nub}[4]{\nu #1{:}#2 := #3.\,#4}
\newcommand{\ext}{\mathcal{E}}
\newcommand{\extt}{\mathcal{T}}
\newcommand{\extk}{\mathcal{K}}
\newcommand{\extc}{\mathcal{C}}
\newcommand{\coerce}[3]{\mathsf{coerce}\,#1\,#2\,#3}
\newcommand{\islarge}{\mathsf{large}}
\newcommand{\nf}{\mathsf{nf}}
\newcommand{\dyntoken}{\mathsf{dyn}_0}
\newcommand{\lint}{\ell_{\mathsf{int}}}
\newcommand{\reduces}{\longrightarrow}
\newcommand{\reducesstar}{\longrightarrow^{*}}
\newcommand{\conv}{\simeq}
\newcommand{\sepsim}{\mathrel{\lesssim}}
\newcommand{\hastype}[3]{#1 \vdash #2 : #3}
\newcommand{\wf}[1]{#1\ \mathsf{wf}}

%% --- listing style for CoC surface syntax ----------------------------------
\definecolor{kw}{rgb}{0.15,0.15,0.6}
\definecolor{cmt}{rgb}{0.3,0.5,0.3}
\lstdefinelanguage{coc}{
  morekeywords={fun,forall,Set,Prop,Kind,Axiom,Infer,Check,Extract,let,in},
  sensitive=true,
  morecomment=[s]{(*}{*)},
  moredelim=[s][\color{cmt}]{(*}{*)}
}
\lstset{
  language=coc,
  basicstyle=\ttfamily\small,
  keywordstyle=\color{kw}\bfseries,
  commentstyle=\color{cmt}\itshape,
  columns=fullflexible,
  keepspaces=true,
  showstringspaces=false,
  breaklines=true,
  xleftmargin=1em
}

\begin{document}

\title{Extracting Blame}
\subtitle{Verified Optimistic Extraction from CoC to System~F$\omega$ + Blame}

\author{Joomy Korkut}
\email{jkorkut@bloomberg.net}
\orcid{0000-0001-6784-7108}
\affiliation{%
  \institution{Bloomberg}
  \city{New York}
  \state{NY}
  \country{USA}
}

\renewcommand{\shortauthors}{Korkut}

\begin{abstract}
Program extraction turns a proof-carrying, dependently typed development into
executable code. Standard Rocq extraction erases dependent structure into
ML-like programs (OCaml, Haskell, or Scheme) whose target types cannot express
the original dependent invariants; verified erasure and compilation projects
reduce the trusted base of that pipeline, but they do not make the lost
dependent structure explicit in the target type system. We ask a complementary
question: how much of a dependently typed program's structure can be preserved
if we extract into a \emph{typed} target, and can the translation be proved
correct?

We present a mechanized extraction and metatheory for the Calculus of
Constructions (CoC) targeting System~F$\omega$ extended with a dynamic type and
blame tracking (``System~F$\omega$ + blame''), a polymorphic blame calculus in the
tradition of Ahmed, Findler, Siek, and Wadler. The target is expressive enough
to keep the parametric-polymorphic skeleton of CoC programs---type abstractions,
type applications, and type-family structure are preserved---while its dynamic
type~$\dyn$ and casts make target-inexpressible residue explicit rather than
silently dropping it. The extraction is \emph{optimistic}: many term-level dependent indices
are erased from target types, type-level indices are kept, and only translation
cases outside the syntactic F$\omega$-skeletal fragment characterized by our
$\dyn$-freedom theorem---most visibly values used as types---introduce $\dyn$;
when such residue is forced in term position, extraction uses a reserved internal
blame label. (The current extractor does not implement Coq-style proof erasure;
proof terms are treated as ordinary residual terms.)

Everything is formalized in Rocq with no admitted lemmas and no project-local
axioms: the target calculus, well-typedness of extraction, the source-reduction
simulation theorem, and derivation independence are all machine-checked, with the
only assumption anywhere being the standard-library \texttt{eq\_rect\_eq}
introduced by dependent case analysis (\S\ref{sec:mech}). We prove that the
extraction of a well-typed source term is well typed in the current permissive
target typing judgment, a one-step simulation stated with a syntactic simulation
relation, derivation
independence, a Jack-of-All-Trades instantiation lemma, and---crucially---%
\emph{external blame freedom}: no programmer-supplied cast boundary can ever
fire; the reserved internal labels are deliberately outside the theorem's
scope.
To make the
formalization concrete we build a REPL that type-checks CoC terms and prints
their extractions, with a verified de~Bruijn-to-named printer, and we exhibit
eleven worked examples---bounds-checked indexing, propositional equality and
transport, dependent pairs, type-safe \texttt{printf}, height-indexed trees,
session types, an intrinsically typed interpreter, dimension-indexed matrices,
units of measure, a Tarski universe, and higher-kinded functors/monads---each
illustrating what survives extraction and where the dynamic type earns its
place. To our knowledge this is
the first proof-assistant mechanization of a polymorphic blame calculus with
F$\omega$-style higher-kinded type syntax, operational determinism, type-level
definitional equality, and progress, and the first verified extraction from a
dependently typed source into a gradually typed target calculus.  It is not the
first extraction from CoC to an F$\omega$-like language---Paulin-Mohring
established that lineage~\cite{paulinmohring1989fomega}---but rather the first
to target a gradual/blame calculus with mechanized extraction and blame-safety
metatheory.
\end{abstract}

\begin{CCSXML}
<ccs2012>
<concept>
<concept_id>10003752.10003790.10011740</concept_id>
<concept_desc>Theory of computation~Type theory</concept_desc>
<concept_significance>500</concept_significance>
</concept>
<concept>
<concept_id>10011007.10011006.10011008.10011009.10011012</concept_id>
<concept_desc>Software and its engineering~Functional languages</concept_desc>
<concept_significance>300</concept_significance>
</concept>
</ccs2012>
\end{CCSXML}
\ccsdesc[500]{Theory of computation~Type theory}
\ccsdesc[300]{Software and its engineering~Functional languages}

\keywords{program extraction, gradual typing, blame, dependent types,
Calculus of Constructions, System F$\omega$, mechanized metatheory, Rocq}

\maketitle

%% ==========================================================================
\section{Introduction}
%% ==========================================================================

Dependently typed languages such as Rocq (formerly Coq), Agda, Lean, and Idris
let a programmer state and prove precise properties of code. To \emph{run} that
code, one erases the static, non-computational parts and emits an executable
residue. Rocq's extraction \cite{letouzey2003extraction,letouzey2004these} is
the canonical example: it maps Gallina to OCaml, Haskell, or Scheme, erasing the
proof/type structure that the target languages cannot represent directly. This is enormously useful, but it makes two concessions.
First, the \emph{target type system cannot express the source's dependent
invariants}: a well-typed dependent program becomes an ML term whose safety is
no longer captured by the target's types. Second, the extractor is part of the
\emph{trusted computing base}: its correctness is argued informally, not proved
inside the proof assistant. (Recent work \cite{forster2024verified} verifies
extraction from Coq to OCaml, reducing this trust; our work is complementary,
studying extraction into a typed target that makes lost dependency explicit.)

The idea of extracting a non-dependent polymorphic program from a CoC
development goes back to Paulin-Mohring's extraction of F$\omega$ programs from
proofs in the Calculus of Constructions~\cite{paulinmohring1989fomega}.  That
work uses a realizability interpretation and a distinction between informative
and non-informative propositions to remove logical content.  We return to the
same high-level source/target tension, but with a different target and theorem:
rather than extracting only the F$\omega$-realizable fragment, we extract
optimistically into F$\omega$ extended with $\dyn$, casts, blame, and
$\nu$-style sealing, so that target-inexpressible residue is explicit and
covered by an external-blame-freedom theorem.

This paper investigates this design point. We fix an expressive
\emph{typed} target and ask how much of a dependently typed program's structure
we can \emph{preserve}, and whether the translation can be \emph{verified}. Our
source is the Calculus of Constructions (CoC) \cite{coquand1988coc}, presented
as a pure type system and mechanized following Barras' \emph{Coq-in-Coq}
development \cite{barras1999,barras1997coqincoq}. Our target is System~F$\omega$
extended with a dynamic type $\dyn$, type-directed casts, and blame tracking:
a \emph{polymorphic blame calculus} in the line of Wadler and Findler
\cite{wadler2009blame} and Ahmed, Findler, Siek, and Wadler
\cite{ahmed2011blame}. We call it \emph{System~F$\omega$ + blame}.

\paragraph{Why a gradual target?}
F$\omega$ preserves a large parametric-polymorphic skeleton: type
abstraction, type application, and higher-kinded type families
\cite{reynolds1983types}. CoC has strictly more---full dependency, in which
types may mention terms. The mismatch is fundamental for a direct,
type-preserving translation into ordinary F$\omega$: arbitrary term dependency is
not part of the target's type formation. A gradual target resolves the tension without
falling back to a fully untyped image: we keep everything F$\omega$ can express
and route the genuine residue of dependency through the dynamic type $\dyn$,
with casts marking the boundary and blame labels attributing responsibility if a
boundary ever fails. The result is a translation that is \emph{optimistic}: it
stays maximally static, becoming dynamic only where it must.

Concretely, three things happen at the boundary (\S\ref{sec:overview}
illustrates all three):
\begin{itemize}
\item \textbf{Term-indexed dependencies are erased from target types.}
A length-indexed vector \texttt{Vec A n} becomes the plain polymorphic type
\texttt{Vec A}: the target type no longer records the relation between the
vector and the term index~$n$.  The term~$n$ itself is not necessarily erased:
if the source program passes it to a function or eliminator, it remains an
ordinary target term argument.  What disappears is the dependent typing
invariant, not every runtime occurrence of the index.
\item \textbf{Type indices are kept.} A type family applied to a type argument
extracts to a type application, preserving parametricity and the family's
higher-kinded structure.
\item \textbf{Values used as types become dynamic residue.}
A source term used where the target expects a type, such as the decoding
function of a Tarski universe, has no direct F$\omega$ type-level image.
In target types this residue becomes $\dyn$; in term position, extraction
uses a reserved internal blame term cast through $\dyn$. It is well typed at
every type and simulates anything, but it is a run-time error by design: it
marks exactly where the source outran the target.
\end{itemize}

\paragraph{What is new.}
The novelty is not that CoC contains an F$\omega$-extractable fragment:
Paulin-Mohring established that lineage~\cite{paulinmohring1989fomega}.  Nor is
the novelty verified erasure to executable code: MetaCoq/MetaRocq, CertiCoq,
\OE{}uf, and Forster et al.'s verified extraction address that
space~\cite{sozeau2020metacoq,anand2017certicoq,mullen2018oeuf,forster2024verified}.
The novelty is the combination of three choices: (1)~an explicit polymorphic
blame target with F$\omega$-style type computation, $\dyn$, casts, and
$\nu$-sealing; (2)~an optimistic extraction that preserves the
F$\omega$-expressible skeleton and marks the rest with internal blame; and
(3)~a Rocq mechanization proving well-typed extraction, source-reduction
simulation, derivation independence, external-blame freedom, operational
determinism, and type-level confluence.

\paragraph{Contributions.}
\begin{itemize}
\item We mechanize an F$\omega$-style \emph{polymorphic blame calculus}
(\S\ref{sec:target}): syntax, a small-step semantics proved \emph{deterministic},
a kinding relation, type-level $\beta$/definitional equality with
\emph{confluence} and \emph{Church--Rosser}, a typing conversion rule, a typing
judgment, subtyping, a simulation relation, and the Blame Theorem.  We further
develop the target's structural metatheory---term and type weakening, term and
type substitution, inversion through conversion, and canonical forms---all in
Rocq with no admitted lemmas or project-local axioms; the assumption set is
audited separately (\S\ref{sec:mech}).  This is a
target variant inspired by Blame for All~\cite{ahmed2011blame}, with $\dyn$
restricted to kind $\kstar$ and with a permissive term typing judgment; we
prove progress for this judgment, but full kind-regular progress/preservation
remains future work (\S\ref{sec:limitations}).

\item We define an \emph{optimistic, conversion-invariant} extraction from CoC
to this target (\S\ref{sec:extraction}).  The design returns to
Paulin-Mohring's CoC-to-F$\omega$ insight~\cite{paulinmohring1989fomega} but
adds an explicit gradual boundary: the F$\omega$-expressible skeleton is
preserved, while target-inexpressible type-level residue is represented through
$\dyn$ and internal blame.  The classification of binders as type- or term-level
uses a semantic, reduction-stable predicate ($\islarge$)---not a syntactic
approximation---and types are normalized before translation, so the translation
respects definitional equality. The extraction recurses directly on \emph{typing
derivations}, which required moving CoC's entire reduction, conversion, and
typing infrastructure from \texttt{Prop} to \texttt{Type}.

\item We prove the metatheory (\S\ref{sec:metatheory}): well-typedness
preservation of the extraction; one-step and multi-step \emph{simulation}
against the target's (syntactic) simulation relation; independence of the
extraction from the choice of typing derivation; a Jack-of-All-Trades lemma
relating any type instantiation to the $\dyn$-instantiation up to simulation;
\emph{external blame freedom}; and an \emph{optimism} result showing the
source-type skeleton fragment extracts without $\dyn$.

\item We build a REPL (\S\ref{sec:repl}) that type-checks CoC terms and prints
their F$\omega$ + blame extractions, with a \emph{verified} de~Bruijn-to-named
renderer whose name recovery is capture-free by construction.

\item We present eleven worked examples (\S\ref{sec:overview},
\S\ref{sec:examples}) that exercise the design: they are the kind of programs
one actually writes with dependent types, and they make visible what the
extraction keeps, erases, and dynamizes.
\end{itemize}

The target calculus, well-typedness of extraction, the source-reduction
simulation theorem, and derivation independence all contain no admitted lemmas
or project-local axioms; the sole assumption in the entire development is the
standard-library \texttt{eq\_rect\_eq} (\S\ref{sec:mech}). In particular, the
type-substitution/extraction commutation up to F$\omega$ definitional equality
used by the application rule's conversion step is fully proved.
We stress up front an
honest scope limitation (\S\ref{sec:limitations}): the extraction is defined as
a total mathematical function on derivations, but it invokes normalization and a
type-checking decision procedure, so it is a metatheoretic object rather than an
efficient compiler pass; and our simulation is stated with the target's
simulation relation, whose bridge to full observational approximation we leave
open. Our aim is to chart the design space of \emph{typed, verified} extraction,
not to ship a production back end.

\paragraph{Scope.}
We study a single design point: verified, typed, optimistic extraction from a
pure-PTS source to a polymorphic blame target.  We do not address proof erasure,
inductive families, universe polymorphism, or compilation to executable code.
\S\ref{sec:limitations} catalogues the specific gaps; each marks a possible
extension, not a limitation of the results we do prove.

%% ==========================================================================
\section{Overview by Example}\label{sec:overview}
%% ==========================================================================

We begin with the phenomenology. Throughout, families and their eliminators are
introduced as \texttt{Axiom}s---the honest way to model inductive families in a
bare pure type system---and functions are ordinary CoC terms whose extraction is
governed by our verified translation. Each \texttt{Extract} line below is
output of our REPL (\S\ref{sec:repl}); type abstractions print as $\Lambda$,
type applications as $[\cdot]$, casts as $\langle A \Rightarrow B\rangle^{\ell}$,
and the dynamic type as \texttt{?}. We elide the pretty-printer's redundant
parentheses for readability; nothing else is edited.

\paragraph{Term-indexed dependencies disappear from target types.}
A bounds-checked lookup indexes a length-indexed vector by a provably in-range
position:
\begin{lstlisting}
Axiom Vec : Set -> Nat -> Set.
Axiom Fin : Nat -> Set.
Axiom vnth : forall (A : Set) (n : Nat), Vec A n -> Fin n -> A.

Extract fun (A : Set) (n : Nat) (xs : Vec A n) (i : Fin n) => vnth A n xs i.
(* ==>  LA:*. \n:Nat. \xs:Vec A. \i:Fin. (vnth [A] n xs i) *)
\end{lstlisting}
The length~$n$ is a term; it appears in the source types \texttt{Vec A n} and
\texttt{Fin n} but is dropped from the extracted types \texttt{Vec A} and
\texttt{Fin}. The term arguments \texttt{n} and \texttt{i} themselves remain
because \texttt{vnth} uses them as ordinary data.  What the target no longer
records is the dependent invariant that \texttt{i} is bounded by~$n$: the
typing relationship is what disappears, not every runtime occurrence of the
index.

\paragraph{Equality indices vanish from types; proof terms remain.}
Propositional equality lives in \texttt{Prop}. The current extractor does not
implement Coq-style proof erasure, but it erases the equality endpoints from
target types: \texttt{Eq A x y} becomes \texttt{Eq A}. Proof arguments and
eliminators such as \texttt{eq\_rect} remain as ordinary residual terms:
\begin{lstlisting}
Extract fun (A : Set) (x y : A) (h : Eq A x y) =>
  eq_rect A x (fun (z : A) => Eq A z x) (refl A x) y h.
(* ==>  LA:*. \x:A. \y:A. \h:Eq A. (eq_rect [A] x [Eq A] (refl [A] x) y h) *)
\end{lstlisting}
In a dependently typed language, \emph{transport} (\texttt{eq\_rect}) is how one
crosses between two provably equal types; it is the static ancestor of a cast.
This example demonstrates index erasure from types, not proof erasure.
A future optimization could erase proof terms more aggressively.

\paragraph{Dependent arities collapse.}
Type-safe \texttt{printf} indexes a format by the type of the variadic function
it denotes; the number and types of arguments are computed from the format
value. The dependency is entirely static:
\begin{lstlisting}
Extract sprintf (Str -> Nat -> Str) (fstr (Nat -> Str) (fint Str fstop)).
(* ==>  (sprintf [Str -> Nat -> Str] (fstr [Nat -> Str] (fint [Str] fstop))) *)
\end{lstlisting}
Here the computed arity \texttt{Str -> Nat -> Str} is a perfectly ordinary
F$\omega$ arrow type, kept as a type argument; the dependent arity leaves no
dynamic checks and is represented as an ordinary F$\omega$ type argument.

\paragraph{Where the dynamic type earns its keep.}
The examples above stay within the target's static skeleton: their dependent
indices are erased from target types without introducing~$\dyn$. The exception
is a value used \emph{as a
type}. A Tarski universe has codes \texttt{U} and a decoder \texttt{El : U ->
Set}; the function \texttt{fun c => El c} maps data to types:
\begin{lstlisting}
Extract fun (c : U) => El c.
(* ==>  \c:U. ((<? => U -> ?>^0+ (<? => ?>^0+ blame(0+))) c) *)
\end{lstlisting}
This is the boundary. \texttt{El} has no F$\omega$ image, so the extraction
emits an internal blame (label~$0$), cast through the dynamic type~$\dyn$, and
applies it to the code. This is not a bug; it is the honest image of ``a value
in type position,'' and our blame-freedom theorem
(\S\ref{sec:metatheory}) proves that no external label can fire---only the
reserved internal labels, never a boundary the programmer wrote.

\medskip
The lesson of these examples, and of the eleven in \S\ref{sec:examples}, is that a
gradual target lets a \emph{single} translation be both faithful and total: it
preserves the polymorphic skeleton, removes many dependent indices from target
types, and localizes the irreducible residue of dependency to a well-understood,
provably contained corner.

%% ==========================================================================
\section{The Source: the Calculus of Constructions}\label{sec:source}
%% ==========================================================================

Our source is the Calculus of Constructions, presented as a pure type system
(PTS) \cite{barendregt1991lambda} with three sorts. We reuse Barras'
\emph{Coq-in-Coq} formalization \cite{barras1999} for terms, de~Bruijn
substitution, reduction, confluence, strong normalization, and decidable
type-checking, adapting it as described in \S\ref{sec:mech}.

\paragraph{Syntax and sorts.}
Terms are the usual PTS terms over de~Bruijn indices:
\[
t, T ::= s \mid n \mid \lambda T.\,t \mid t\ t \mid \Pi T.\,T,
\qquad
s ::= \Prop \mid \Set \mid \Kind.
\]
The sort axioms are $\Prop : \Kind$ and $\Set : \Kind$; the top sort $\Kind$ is
untypable, which is what averts Girard's paradox \cite{girard1972}. The product
rule is unrestricted over the typable sorts: for all $s_1, s_2$,
\[
\inferrule{\hastype{\Gamma}{T}{s_1} \\ \hastype{\Gamma, T}{U}{s_2}}
          {\hastype{\Gamma}{\Pi T.\,U}{s_2}}
\]
so products are impredicative in both $\Prop$ and $\Set$. Figure~\ref{fig:coc}
gives the typing rules in full. Conversion is definitional: the rule
\textsc{Conv} lets a term inhabit any type convertible ($\conv$, the
symmetric-transitive closure of $\beta$-reduction) to its inferred one.

\begin{figure}
\small
\begin{mathpar}
\inferrule*[right=\textsc{Prop}]{\wf{\Gamma}}{\hastype{\Gamma}{\Prop}{\Kind}}
\and
\inferrule*[right=\textsc{Set}]{\wf{\Gamma}}{\hastype{\Gamma}{\Set}{\Kind}}
\and
\inferrule*[right=\textsc{Var}]{\wf{\Gamma} \\ \Gamma(n) = T}{\hastype{\Gamma}{n}{T}}
\\
\inferrule*[right=\textsc{Lam}]
  {\hastype{\Gamma}{T}{s_1} \\ \hastype{\Gamma, T}{U}{s_2} \\ \hastype{\Gamma, T}{M}{U}}
  {\hastype{\Gamma}{\lambda T.\,M}{\Pi T.\,U}}
\and
\inferrule*[right=\textsc{App}]
  {\hastype{\Gamma}{u}{\Pi V.\,U_r} \\ \hastype{\Gamma}{v}{V}}
  {\hastype{\Gamma}{u\ v}{U_r[v]}}
\\
\inferrule*[right=\textsc{Prod}]
  {\hastype{\Gamma}{T}{s_1} \\ \hastype{\Gamma, T}{U}{s_2}}
  {\hastype{\Gamma}{\Pi T.\,U}{s_2}}
\and
\inferrule*[right=\textsc{Conv}]
  {\hastype{\Gamma}{t}{U} \\ U \conv V \\ \hastype{\Gamma}{V}{s}}
  {\hastype{\Gamma}{t}{V}}
\end{mathpar}
\caption{Typing for the source Calculus of Constructions (de~Bruijn; $\Gamma(n)$
denotes the appropriately lifted $n$-th binding). Well-formedness $\wf{\Gamma}$
threads through the rules.}
\label{fig:coc}
\Description{Typing rules for the source Calculus of Constructions.}
\end{figure}

\paragraph{The large/small distinction.}
Extraction hinges on separating \emph{type-level} from \emph{term-level}
objects. A source object is \emph{large} when its type is the top sort $\Kind$:
\[
\islarge(\Gamma, T) \;\triangleq\; \hastype{\Gamma}{T}{\Kind}.
\]
Large things are kinds---$\Set$, $\Prop$, $\Nat \to \Set$, and so on---which
classify types; they must become F$\omega$ type-level binders and quantifiers.
Small things (sort $\Prop$ or $\Set$) classify terms and become term-level
binders. The crucial property, which we prove (\S\ref{sec:metatheory}), is that
$\islarge$ is \emph{stable under reduction}: subject reduction guarantees that
if $T$ has sort $\Kind$ and $T \reduces T'$, then $T'$ has sort $\Kind$. This is
what makes every classification decision in the extraction respect definitional
equality. A purely syntactic approximation ($\mathsf{classifier}$, used only in
positions already known to be normal) is sound but not reduction-stable, and we
are careful never to base a conversion-sensitive decision on it.

%% ==========================================================================
\section{The Target: an F$\omega$-Style Blame Calculus}\label{sec:target}
%% ==========================================================================

The target follows the polymorphic blame calculus of Ahmed et
al.~\cite{ahmed2011blame}, lifted from System~F to F$\omega$-style
higher-kinded type syntax and extended with a dynamic type $\dyn$,
type-directed casts, and blame. We give a machine-checked
formalization of this F$\omega$-style blame variant; the full development lives in the
\texttt{BlameFOmega} directory. The type language has the essential F$\omega$ structure: it has
type operators ($\Lambda$/juxtaposition), a kinding relation, a type-level
$\beta$ rule with its definitional-equality closure, and a typing conversion
rule (\S\ref{sec:limitations} details what is and is not yet mechanized about
the term-level type system). We are precise about scope: the term typing
judgment used by our extraction theorem does not yet embed kinding premises,
so it is not kind-regular. We do prove progress for this judgment; a target
subject-reduction theorem remains future work (\S\ref{sec:limitations}).

\paragraph{Kinds, types, terms.}
\begin{align*}
K &::= \kstar \mid K \karr K & \text{(kinds)}\\
A, B &::= \alpha \mid A \to B \mid \forall \alpha{:}K.\,A \mid \Lambda \alpha{:}K.\,A \mid A\ B \mid \dyn & \text{(types)}\\
e &::= x \mid \lambda x{:}A.\,e \mid e\ e \mid \Lambda \alpha{:}K.\,e \mid e\,[A] \\
  &\quad \mid \cast{e}{A}{B}{p} \mid \gnd{e}{G} \mid \isgnd{e}{G} \mid \blame p \mid \nub{\alpha}{K}{A}{e} & \text{(terms)}
\end{align*}
Types include the polymorphic quantifier $\forall$, the type-operator
abstraction/application $\Lambda$/juxtaposition (needed to preserve CoC's
higher-kinded type families), and the dynamic type $\dyn$. The dynamic type is a
proper type of kind $\kstar$; we do not include a kind-indexed family
$\dyn_K$. Terms extend F$\omega$ with: a cast $\cast{e}{A}{B}{p}$ that coerces
$e$ from $A$ to $B$ under blame label $p$; ground injection $\gnd{e}{G}$ and the
ground test $\isgnd{e}{G}$, which mediate between $\dyn$ and \emph{ground types}
$G ::= \dyn \to \dyn \mid N$, where $N ::= \alpha \mid N\,A$ ranges over
\emph{neutral} types (a type variable or a type variable applied to further
types)---the F$\omega$ generalization of Blame-for-All's type-variable ground
tag, needed once type applications, not just variables, can appear at the
$\dyn$ boundary; a blame term $\blame p$; and a
\emph{type-sealing} binder $\nub{\alpha}{K}{A}{e}$ that introduces a fresh
abstract type variable $\alpha$ of kind $K$ whose hidden implementation is $A$,
used to model polymorphic instantiation dynamically (it is the reduct of a type
application $\Lambda$-elimination). A blame label $p$ pairs an identifier with a
polarity. Identifier~$0$ is reserved for \emph{internal} failures---both the
extraction-internal casts and $\nu$-tampering---and identifier~$1$ for
$\mathsf{is\_gnd}$-tampering; all identifiers $\ge 2$ are programmer-supplied
(\emph{external}) boundaries. (The extraction and $\nu$-tamper labels share
id~$0$; this is harmless for the external-blame theorem, which quantifies over
$\ge 2$, but we note it explicitly.)

\paragraph{Executable compatibility.}
Unlike Blame for All's System-F \emph{consistency} relation, our $\mathsf{compat}(A,B)$
is an \emph{executable} judgment: every constructor corresponds directly to a
cast-reduction rule, so a well-typed cast always has an operational step, with no
catch-all ``stuck cast'' case (except for identity casts, which step without
inspecting their annotation, each constructor decomposes the cast on a canonical
head). This is one of the places where the F$\omega$ extension is not routine: in
System~F, Blame-for-All compatibility can be read as both consistency and a cast
strategy, but once type-level functions and neutral type-family applications are
present, castability must be made executable---every compatibility constructor
becomes a cast-elaboration rule. $\mathsf{compat\_refl}$
licenses the identity cast; $\mathsf{compat\_arrow}$ (contravariant domain,
covariant codomain) licenses the structural WRAP step; $\mathsf{compat\_all}$
licenses the analogous structural step under a matching $\forall$-binder;
$\mathsf{compat\_generalize}$ and $\mathsf{compat\_instantiate}$ license casting
into and out of a $\forall$; $\mathsf{compat\_to\_dyn}$ licenses injection at a
ground tag, itself requiring $\mathsf{compat}(A,G)$ so injection only ever
manufactures a cast the rest of the judgment already accepts; and
$\mathsf{compat\_from\_dyn}$ licenses projecting a ground value back out of
$\dyn$. The judgment is decidable ($\mathsf{compat\_dec}$), which is what makes
$\coerce{\cdot}{\cdot}{\cdot}$ (\S\ref{sec:extraction}) a total function.

\paragraph{Semantics, typing, and safety.}
We define a small-step reduction $[e \reduces e']$ (with values, $\beta$/type-$\beta$
reduction, cast reductions that split structural casts and fire blame on ground
mismatch, and $\nu$-sealing), a kinding relation
$\mathsf{wf\_typ}(\Delta, A, K)$, and a typing judgment
$\hastype{\Delta}{e}{A}$ over contexts $\Delta$ binding term variables (at a
type), type variables (at a kind), and type-definition bindings introduced by
$\nu$-sealing. The dynamic type is inhabited by ground
injections, and casts are well typed between compatible types. The typing judgment
we use for the extraction theorem is deliberately permissive: it does not embed
$\mathsf{wf\_typ}$ premises, so it certifies term-formation rather than full
kind-regularity; strengthening it to a kind-regular target with a
$\mathsf{typing\_regular}$ lemma is future work (\S\ref{sec:limitations}). We
port the subtyping and \emph{safety} apparatus of \cite{ahmed2011blame}: a
predicate $\mathsf{safe}(p, e)$ ensures every cast in $e$ that could blame $p$
carries the right positive/negative subtyping evidence, and safety is preserved
by reduction. From this we obtain the central guarantee.

\begin{theorem}[Blame Theorem, following Thm.~14 of \cite{ahmed2011blame}, mechanized]
\label{thm:blame}
If $\mathrm{lbl\_id}(p) \ge 2$ and $\mathsf{safe}(p, e)$, then
$\neg\,[e \reducesstar \blame p]$.
\end{theorem}

\paragraph{Simulation.}
To state extraction correctness we use a \emph{syntactic simulation} relation
$s \sepsim t$ (Rocq: \texttt{sim}), which relates terms differing by replacing a
type argument with $\dyn$, by inserting or stripping casts, by unfolding
$\beta$/type-$\beta$ redexes, and---critically---by the rule
$\blame q \sepsim t$ for \emph{any} $t$ ($\mathsf{sim\_blame}$): a blame term
is related by $\sepsim$ to every term. Its reflexive-transitive closure is $\sepsim^{*}$
(\texttt{sim\_star}). We use $\sepsim$ as a \emph{syntactic simulation relation}
to state extraction correctness. We define a contextual
approximation (\texttt{approximates}) in the development, but the bridge
$s \sepsim t \Rightarrow s$ approximates $t$ is
\emph{not} mechanized and remains open---indeed, for the current untyped
$\sepsim$ (which includes rules such as
$\mathsf{sim\_left\_tabs}\colon s \sepsim t \Rightarrow \Lambda K.\,s \sepsim t$)
it is false without a typing restriction, since $\Lambda K.\,s$ and $s$ are
distinguished by the context $[\cdot]\,[\dyn]$. We therefore treat every
correctness statement below as a statement about $\sepsim^{*}$, and discuss the
gap to observational approximation in \S\ref{sec:limitations}. The key algebraic
fact we reuse is Jack-of-All-Trades (\S\ref{sec:metatheory}): a type application
at a concrete type is $\sepsim$-related to the same application at $\dyn$.

%% ==========================================================================
\section{The Extraction}\label{sec:extraction}
%% ==========================================================================

The extraction comprises three mutually supporting pieces: a translation of
kinds, a translation of types, and a translation of terms driven by typing
derivations. All three respect definitional equality: types are normalized
before translation, and binder classification uses the reduction-stable
$\islarge$.

\paragraph{Kinds.}
A large source type (a kind) becomes a target kind by turning classifier domains
into kind arrows and dropping term domains:
\[
\extk(s) = \kstar, \quad
\extk(\Pi T.\,U) = \begin{cases}
  \extk(T) \karr \extk(U) & \text{if $T$ is a classifier} \\
  \extk(U) & \text{otherwise.}
\end{cases}
\]
Thus $\Set \to \Nat \to \Set$ becomes $\kstar \karr \kstar$: the type argument
$\Set$ contributes a $\kstar$, the term argument $\Nat$ is dropped. This is
exactly the kind a vector-like family needs.

\paragraph{Types.}
Type translation is $\extt(\Gamma, T) = \extt_L(\Gamma, \nf(T))$: normalize with
$\nf$ (justified by strong normalization), then translate the normal form with
$\extt_L$ (Figure~\ref{fig:extt}). A sort becomes $\dyn$; a variable becomes a
type variable if its binding is large, else $\dyn$; a product becomes a
$\forall$ (large domain) or an arrow (small domain); a type-operator
$\lambda$ becomes $\Lambda$ (large domain) or is erased (small domain); and an
application becomes a type application when both sides are type-level, drops a
term argument when only the head is type-level, and otherwise becomes $\dyn$.
The auxiliary $\mathsf{typeexpr}$ decides whether a term denotes a type-level
object. Because classification here consults $\islarge$ through the context and
because the argument is already normal, $\extt$ is invariant under conversion of
both the type and the context.

\begin{figure}
\small
\begin{align*}
\extt_L(\Gamma, s) &= \dyn \\
\extt_L(\Gamma, n) &= \begin{cases}
   \alpha_{\,\mathrm{idx}(\Gamma,n)} & \text{if $\islarge$ at $n$} \\
   \dyn & \text{otherwise}
 \end{cases}\\
\extt_L(\Gamma, \Pi T.\,U) &= \begin{cases}
   \forall{:}\extk(T).\,\extt_L(\Gamma{,}T, U) & \text{$T$ classifier}\\
   \extt_L(\Gamma, T) \to \extt_L(\Gamma{,}T, U) & \text{otherwise}
 \end{cases}\\
\extt_L(\Gamma, \lambda T.\,M) &= \begin{cases}
   \Lambda{:}\extk(T).\,\extt_L(\Gamma{,}T, M) & \text{$T$ classifier}\\
   \extt_L(\Gamma{,}T, M) & \text{otherwise}
 \end{cases}\\
\extt_L(\Gamma, u\ v) &= \begin{cases}
   \extt_L(\Gamma,u)\ \extt_L(\Gamma,v) & \text{$u,v$ type-level}\\
   \extt_L(\Gamma, u) & \text{$u$ type-level, $v$ not}\\
   \dyn & \text{otherwise}
 \end{cases}
\end{align*}
\caption{Type translation $\extt_L$ on normal forms. Term arguments of a
type-level head are dropped; a type-level application is kept; everything else
falls back to $\dyn$.}
\label{fig:extt}
\Description{Translation rules for source types into target types.}
\end{figure}

\paragraph{Terms.}
The term translation $\ext$ recurses on a \emph{typing derivation}
(Figure~\ref{fig:ext}). Two helpers absorb the impedance mismatch. The
\emph{coercion} $\coerce{s}{A}{B}$ is the identity when $A = B$; otherwise, if
the executable compatibility judgment $\mathsf{compat}(A,B)$ holds (decidable,
\S\ref{sec:target}), a single direct cast; and otherwise reserved internal
blame, an honest representation of a target-inexpressible coercion:
\[
\coerce{s}{A}{B} =
  \begin{cases}
    s & A = B \\
    \cast{s}{A}{B}{\lint} & A \neq B,\ \mathsf{compat}(A,B) \\
    \blame \lint & \text{otherwise.}
  \end{cases}
\]
The \emph{placeholder} $\dyntoken = \gnd{(\lambda{:}\dyn.\,0)}{(\dyn \to \dyn)}$
is an inert, closed, well-typed inhabitant used where a source sort or product
sits in term position.

The interesting cases are \textsc{Var} and \textsc{App}. A variable whose type
is large is a \emph{type variable used as a term}---the one construct with no
sound F$\omega$ image---and extracts to $\coerce{(\blame \lint)}{\dyn}{\cdot}$:
an \emph{internal} blame. A small variable extracts to a coerced term variable.
An application whose function domain is large is a type application ($\Lambda$
elimination) and extracts to a \emph{bare} type application $e\,[A]$: its
natural target type $\extt(V{:}\Gamma, U_r)[\extt(\Gamma,v)/0]$ and the expected
result type $\extt(\Gamma, U_r[v])$ are equal only up to F$\omega$ definitional
equality $\mathsf{ty\_equiv}$ (\S\ref{sec:target}), not on the nose, so this case
relies on the target conversion rule \textsc{T-Conv} rather than a runtime
coercion (see ``extraction does not commute with substitution syntactically''
below). A small application extracts to a coerced $e\ e'$. Conversions extract to
coercions between the extracted source and target types---which, by
conversion-invariance of $\extt$, are equal, so the coercion is the identity.
Sorts and products in term position extract to $\dyntoken$.

\begin{figure}
\small
\begin{align*}
\ext(\textsc{Prop}) &= \ext(\textsc{Set}) = \dyntoken \\
\ext(\textsc{Var}\ n) &= \begin{cases}
   \coerce{(\blame \lint)}{\dyn}{\extt(\Gamma, T)} & \islarge(\Gamma, T)\\
   \coerce{x_{\mathrm{idx}(\Gamma,n)}}{\extt_{\mathrm{lookup}}(\Gamma, n)}{\extt(\Gamma, T)} & \text{otherwise}
 \end{cases}\\
\ext(\textsc{Lam}) &= \begin{cases}
   \Lambda{:}\extk(T).\,\ext(M) & \islarge(\Gamma, T)\\
   \lambda{:}\extt(\Gamma, T).\,\ext(M) & \text{otherwise}
 \end{cases}\\
\ext(\textsc{App}) &= \begin{cases}
   \ext(u)\,[\extt(\Gamma, v)] & \islarge(\Gamma, V)\\
   \coerce{(\ext(u)\ \ext(v))}{\cdot}{\cdot} & \text{otherwise}
 \end{cases}\\
\ext(\textsc{Prod}) &= \dyntoken \\
\ext(\textsc{Conv}) &= \coerce{\ext(t)}{\extt(\Gamma, U)}{\extt(\Gamma, V)}
\end{align*}
\caption{Term translation $\ext$, by recursion on the typing derivation. The
large \textsc{App} case emits a bare type application, well typed at its natural
type $\extt(V{:}\Gamma,U_r)[\extt(\Gamma,v)/0]$ and reconciled with the expected
result type $\extt(\Gamma,U_r[v])$ by the target conversion rule \textsc{T-Conv}
(the two are $\mathsf{ty\_equiv}$-equal but not syntactically equal). The elided
coercion targets in the small cases are the substituted vs.\ actual result types;
there they are equal by conversion-invariance of $\extt$, so the coercion is the
identity. The elided source types $T, V, U$ are those of the corresponding
derivation nodes.}
\label{fig:ext}
\Description{Translation rules for source typing derivations into target terms.}
\end{figure}

\paragraph{Extraction does not commute with substitution syntactically.}
The large \textsc{App} case rests on a subtle point: type extraction commutes
with the source substitution of the application rule only \emph{up to}
$\mathsf{ty\_equiv}$, never on the nose. Concretely, for $u : \Pi V.\,U_r$ with
$V$ large and argument $v : V$, the two sides
\[
  \extt(V{:}\Gamma, U_r)\,[\,\extt(\Gamma, v)\,/\,0\,]
  \qquad\text{and}\qquad
  \extt(\Gamma,\, U_r[v])
\]
are $\mathsf{ty\_equiv}$-equal but syntactically distinct, for two independent
reasons. \emph{(i) Target $\beta$.} Plugging a large argument can place a
type-operator abstraction in head position, so the left side may contain a
target redex $(\Lambda \alpha{:}K.\,A)\ B$ whose contractum is what appears on the right;
these are related only by the F$\omega$ type-reduction $\mathsf{ty\_step}$.
\emph{(ii) Annotation drift under $\nf$.} Source substitution rewrites binder
annotations up to conversion, and $\extt = \extt_L \circ \nf$ collapses those to
a common normal form; the raw (pre-$\nf$) syntaxes differ. A naive
substitution-commutation lemma phrased with syntactic equality is therefore
\emph{false}. The correct statement (Rocq: \texttt{extract\_typ\_tsubst\_coc\_equiv})
replaces $=$ with $\mathsf{ty\_equiv}$, and its $\beta$ case is discharged by
$\mathsf{ty\_equiv\_beta}$. This is exactly why the extractor emits a bare $e\,[A]$
and defers to \textsc{T-Conv}: the target's F$\omega$ conversion rule
absorbs precisely this definitional equality, so no runtime cast is needed.
The same phenomenon defeats the naive \emph{kind} index of type extraction
(\texttt{extract\_kind\_L} does not commute with the substitution either, since
substituting a type for a variable can flip a domain's syntactic classifier), so
kind-regularity (\texttt{extract\_typ\_wf\_sort}) must thread the target kind
through the CoC kinding derivation rather than read it off the CoC type
syntactically.

\paragraph{Optimism, made precise.}
Two design choices make the extraction \emph{optimistic}---as static as
possible. First, term indices of type families are \emph{dropped}, not
dynamized: \texttt{Vec A n} extracts to the F$\omega$ type application
$\extt(\Gamma,\texttt{Vec})\,\extt(\Gamma,\texttt{A})$, not to something
mentioning $\dyn$. Second, the only sources of $\dyn$ in extracted \emph{types}
are the genuinely irreducible ones in Figure~\ref{fig:extt} (a sort, a
term-level variable, or a term-headed application appearing in type position).
On well-typed closed terms these do not arise from the erasable idioms of
\S\ref{sec:overview}; they arise exactly at the value-as-type boundary, where
the term translation additionally emits an internal $\blame$. We quantify the
optimism formally in Theorem~\ref{thm:dynfree}.

%% ==========================================================================
\section{Metatheory}\label{sec:metatheory}
%% ==========================================================================

All results below are mechanized in Rocq and are admit-free (the only
assumption anywhere is the standard-library \texttt{eq\_rect\_eq};
\S\ref{sec:mech}). We name each
Rocq theorem for cross-reference. Let $\ext(\Gamma \vdash t : T)$ denote the
term extraction of a derivation of $\hastype{\Gamma}{t}{T}$, and
$\extc(\Gamma)$ the extracted context.

\begin{theorem}[Well-typed extraction; \texttt{extract\_well\_typed}]
\label{thm:welltyped-extraction}
For every derivation of $\hastype{\Gamma}{t}{T}$ with $\Gamma$ well-formed,
\[
\hastype{\extc(\Gamma)}{\ext(\Gamma \vdash t : T)}{\extt(\Gamma, T)}.
\]
\end{theorem}
The extraction of a well-typed CoC term is well typed in the current permissive
target typing judgment, at the extracted type. (This is \emph{well-typedness of
the extraction}, not target subject reduction; and the target judgment does not
embed full kinding premises, as discussed in \S\ref{sec:target}.) The proof is by induction on the derivation, using
conversion-invariance of $\extt$ to discharge the coercion cases and the typing
of $\dyntoken$ and $\blame \lint$ at $\dyn$.

\begin{theorem}[Simulation; \texttt{extract\_reduces\_once}, \texttt{extract\_reduces}]
\label{thm:sim}
If $\hastype{\Gamma}{t}{T}$ and $t \reduces v$ (so, by subject reduction,
$\hastype{\Gamma}{v}{T}$), then
\[
\ext(\Gamma \vdash t : T) \;\sepsim^{*}\; \ext(\Gamma \vdash v : T).
\]
The multi-step version follows by transitivity of $\sepsim^{*}$.
\end{theorem}
Source reduction is tracked by target simulation. The proof develops the
expected substitution lemmas at the target: extraction commutes with term
substitution and with type substitution up to $\sepsim$, and $\beta$/type-$\beta$
source redexes map to the corresponding $\mathsf{sim\_beta}$/$\mathsf{sim\_tbeta}$
steps. Weakening lemmas (\texttt{extract\_weaken}$\ast$) handle the shift in
de~Bruijn indices between the two namespaces.

This theorem is deliberately \emph{syntactic}. Because $\sepsim$ contains the
one-sided rule $\blame q \sepsim t$, an extraction that has already fallen to
reserved internal blame on the left simulates the reduct trivially: the relation
records that the target has entered an internal-residue case rather than
asserting behavioral preservation. A behavioral (contextual-approximation)
statement would need a typed reformulation of $\sepsim$ and a CIU/logical-relations
adequacy argument; that bridge is open (\S\ref{sec:limitations}).

\begin{theorem}[Derivation independence; \texttt{extract\_deriv\_indep}]
\label{thm:indep}
Any two derivations $H_1, H_2$ of the same judgment $\hastype{\Gamma}{t}{T}$
satisfy $\ext(\Gamma \vdash_{H_1} t : T) \sepsim^{*} \ext(\Gamma \vdash_{H_2} t :
T)$.
\end{theorem}
Because $\ext$ recurses on derivations, this is what licenses talking about
``the'' extraction of a term: different derivations agree up to simulation. It
follows from a conversion-strengthening lemma
(\texttt{extract\_deriv\_indep\_conv}) applied to reflexivity.

\begin{theorem}[External blame freedom; \texttt{extraction\_blame\_free}]
\label{thm:blamefree}
For every derivation of $\hastype{\Gamma}{t}{T}$ and every label $p$ with
$\mathrm{lbl\_id}(p) \ge 2$,
\[
\neg\,\bigl[\ext(\Gamma \vdash t : T) \reducesstar \blame p\bigr].
\]
\end{theorem}
No extracted term ever blames a programmer-supplied (external) boundary. The
proof combines a safety lemma (\texttt{extracted\_safe}: every cast the
extraction introduces uses the internal label $\lint$, hence is safe for any
external $p$) with the Blame Theorem (Theorem~\ref{thm:blame}). The reserved
internal identifier~$0$ is deliberately \emph{not} covered. Two sources can emit
it: the value-as-type residue the extraction inserts as $\blame\lint$
(\S\ref{sec:extraction}), and---since $\mathsf{nu\_tamper\_label}$ shares
id~$0$---a $\nu$-tamper in the raw target semantics, now generalized
(\S\ref{sec:target}) to fire whenever a sealed variable would be exposed
through \emph{any} neutral ground tag $N$, not only a bare type variable.  Both
are internal target failures, never programmer boundaries; the theorem
excludes exactly this reserved identifier.

\begin{theorem}[Jack-of-All-Trades; \texttt{extraction\_simulates\_any\_instantiation}]
\label{thm:jack}
For any concrete type $C$ and a kind-$\kstar$ quantifier, instantiating an
extracted polymorphic term at $\extt(\Gamma, C)$ is simulated by instantiating
it at $\dyn$:
\[
\cast{(\ext(M)\,[\extt(\Gamma,C)])}{\_}{B}{p}
\;\sepsim\;
\cast{(\ext(M)\,[\dyn])}{\_}{B}{p}.
\]
\end{theorem}
This is the extraction-level instance of the target's Jack-of-All-Trades
simulation lemma (\texttt{jack\_of\_all\_trades}): the concrete instantiation is
$\sepsim$-related to the $\dyn$ instantiation. It is a \emph{syntactic}
simulation principle; turning it into an observational statement about
parametric behavior would require the open $\sepsim$-to-approximation bridge
(\S\ref{sec:limitations}).

\begin{theorem}[Optimism / $\dyn$-freedom; \texttt{extract\_typ\_dyn\_free\_iff}]
\label{thm:dynfree}
On the source-type skeleton fragment $\mathsf{fo}$---the fragment whose
surviving type-level structure is F$\omega$-expressible (variables, arrows, and
type applications) and whose term-level family arguments occur only in positions
erased by type extraction---type extraction avoids $\dyn$ \emph{if and only if}
the normal form lies in $\mathsf{fo}$:
\[
\mathsf{typ\_dyn\_free}(\extt(\Gamma, T))
\;\Longleftrightarrow\;
\mathsf{fo}(\Gamma, \nf(T)).
\]
\end{theorem}
This is an exact characterization: the fragment $\mathsf{fo}$ captures
precisely the source types whose extraction is $\dyn$-free. Note that
$\mathsf{fo}$ includes many indexed-family uses such as \texttt{Vec A n},
because the term index \texttt{n} is erased from the target type
($\extt(\Gamma,\texttt{Vec A n}) = \extt(\Gamma,\texttt{Vec})\,\extt(\Gamma,\texttt{A})$);
it excludes term-headed type computation such as \texttt{El c}, where a term
determines a type.
On the fragment where F$\omega$ is genuinely as expressive as the source, the
translation is not merely optimistic but $\dyn$-free: the target type is fully
static. Outside the fragment $\dyn$ must appear, exactly as
Figure~\ref{fig:extt} prescribes.

\medskip
Taken together: extraction produces well-typed target terms
(\ref{thm:welltyped-extraction}), tracks computation up to $\sepsim^{*}$
(\ref{thm:sim}), is canonical (\ref{thm:indep}), never blames the programmer
(\ref{thm:blamefree}), relates any instantiation to the $\dyn$-instantiation by
$\sepsim$ (\ref{thm:jack}), and is $\dyn$-free exactly on the source-type skeleton fragment
(\ref{thm:dynfree}).

%% ==========================================================================
\section{Mechanization}\label{sec:mech}
%% ==========================================================================

The development is written in Rocq~9.0 and builds with \texttt{dune}, and is
\emph{entirely admit-free}: the target calculus (\texttt{BlameFOmega}),
well-typed extraction (\texttt{extract\_well\_typed}), the source-reduction
simulation (\texttt{extract\_reduces\_once}/\texttt{extract\_reduces},
Theorem~\ref{thm:sim}), and the derivation-independence and blame-freedom
theorems all carry complete machine-checked proofs. The extraction
weakening/substitution commutation lemmas the simulation needs---%
\texttt{term\_tlift\_extract\_sim}, \texttt{term\_tsubst\_extract\_sim},
\texttt{extract\_weaken1}, \texttt{extract\_subst\_sim\_gen},
\texttt{extract\_tsubst\_gen}---are the ``$\nf$ commuting with
\texttt{lift}/\texttt{subst}'' inductions, parallel to
\texttt{extract\_typ\_L\_swap}'s treatment of reduction.

The two lemmas discharging the large-\textsc{App} \textsc{T-Conv} premises are
both proved. Kind regularity, \texttt{extract\_typ\_wf\_sort} (type extraction is
well kinded at $\kstar$), is axiom-free via a stronger induction
(\texttt{extract\_typ\_L\_wf\_kind}) that threads the target kind through the CoC
kinding derivation restricted to genuinely type-level subterms
(\texttt{type\_expr}), which is what makes the otherwise-false naive kind index
correct. The conversion premise, \texttt{extract\_typ\_tsubst\_coc\_equiv}
(type extraction commutes with the application rule's source type substitution
\emph{up to} $\mathsf{ty\_equiv}$; \S\ref{sec:extraction}), factors into a raw
structural substitution-commutation lemma
(\texttt{extract\_typ\_L\_large\_subst\_raw}, plain structural induction---%
\texttt{extract\_typ\_L} never fires a redex) and reduction-invariance of
\texttt{extract\_typ\_L} up to $\mathsf{ty\_equiv}$
(\texttt{extract\_typ\_L\_reduces\_nf\_equiv}, a single-step lemma chained over
$W \to^{*} \nf\,W$); the single-step $\beta$ case defers to the raw lemma via
$\mathsf{ty\_equiv\_beta}$. The repository includes
\texttt{Extraction/Assumptions.v}, which audits the axiom footprint of every
headline theorem; the only assumption that appears is the standard library's
\texttt{eq\_rect\_eq} instance introduced by dependent \texttt{destruct} (the
REPL name-recovery lemmas additionally use the standard primitive
string/integer axioms, since they manipulate variable names; the core metatheory
does not). The original Coq-in-Coq base itself remains zero-admit.

\paragraph{Type-valued derivations.}
The single most consequential engineering decision is that $\ext$ recurses on a
typing \emph{derivation}, and a derivation carries computational content only if
the typing judgment lives in \texttt{Type}, not \texttt{Prop}: Rocq forbids
eliminating a \texttt{Prop} inductive into a \texttt{Type}-returning function.
We therefore moved \texttt{has\_type} and \texttt{well\_formed} to \texttt{Type},
which cascaded through the whole reduction/conversion layer: \texttt{reduces},
\texttt{reduces\_once}, \texttt{convertible}, and the parallel-reduction
relations all became \texttt{Type}-valued, with hand-rolled inductives replacing
the standard-library \texttt{clos\_*} closures. Strong normalization, which is
naturally a \texttt{Prop}, is bridged to the \texttt{Type} world by a wrapper
(\texttt{reduces\_once\_prop}). Numerous \texttt{exists}/$\vee$ were changed to
$\Sigma$-types/$+$ across roughly twenty files, all while retaining the original
proofs' structure.

\paragraph{Normalization and decidability inside the function.}
The extraction calls $\nf$ (compute a normal form, total by strong
normalization) and $\islarge$-decision (\texttt{is\_large\_dec}, the type-checking
decision procedure) to classify context bindings in a conversion-invariant way.
These are well-defined total functions, so $\ext$ is a total mathematical
function on derivations; but they make $\ext$ a \emph{metatheoretic} object
rather than a term that reduces under \texttt{reflexivity}
(\S\ref{sec:limitations}). Conversion-invariance is itself a theorem:
$\islarge$ is stable under reduction (\texttt{is\_large\_stable}), and $\extt$ is
invariant under reduction of both its type argument and its context
(\texttt{extract\_typ\_L\_swap}, \texttt{extract\_ctx\_swap}).

\paragraph{Proof effort.}
The \texttt{BlameFOmega} target (syntax, semantics, kinding/typing, subtyping,
safety, the Blame Theorem, and the simulation relation with Jack-of-All-Trades)
is a self-contained mechanization of a polymorphic blame calculus. The
\texttt{Extraction} layer (the translation and \S\ref{sec:metatheory}) is the
largest single file, dominated by the substitution/weakening infrastructure
underlying Theorem~\ref{thm:sim}. The CoC layer is Barras' development with the
\texttt{Type}-valued adaptation.

%% ==========================================================================
\section{A REPL with Verified Name Recovery}\label{sec:repl}
%% ==========================================================================

To make the extraction inspectable we extract the type checker and the
translation to OCaml (via Rocq's extraction) and wrap them in a REPL that reads
CoC surface syntax and prints inferred types and F$\omega$ + blame extractions.
All output in this paper is produced by it. The commands are \texttt{Axiom},
\texttt{Infer}, \texttt{Check}, and \texttt{Extract}; \S\ref{sec:examples}
lists complete input files.

\paragraph{Verified de~Bruijn-to-named printing.}
Metatheory uses de~Bruijn indices, but readable output needs names. Naively
re-inventing names (\texttt{x0}, \texttt{x1}, \dots) is unreadable and, if done
by hand in the printer, unverified and capture-prone. We instead verify the
conversion. On the source side, \texttt{expression\_of\_term\_with\_hints}
converts a de~Bruijn term to a named expression, reusing the programmer's
original binder names as hints; on the target side,
\texttt{fterm\_expression\_of} does the same for F$\omega$ + blame across its two
namespaces (type and term variables). Both are proved to produce output
\emph{equivalent} to the de~Bruijn input (an inductive equivalence relation
carried in the function's result type), and the shared name chooser
\texttt{pick\_name} accepts a hint only when it is fresh for the current context,
so recovered names are \emph{capture-free by construction}---the freshness proof
is part of \texttt{pick\_name}'s type. We further prove the hinted conversion is
alpha-equivalent to the fresh-name conversion (\texttt{..\_alpha}), and that a
named expression determines its de~Bruijn preimage uniquely
(\texttt{..\_unique}). The REPL's printer is thus a thin unverified shell around
verified conversions.

%% ==========================================================================
\section{Worked Examples}\label{sec:examples}
%% ==========================================================================

Table~\ref{tab:examples} summarizes the eleven examples shipped with the
development. Each is a self-contained \texttt{.v} file that type-checks and runs
through the REPL; together they span the idioms one actually writes with
dependent types, and they make the extraction's behavior concrete.

\begin{table}
\small
\caption{The eleven examples. ``Erased'' names the dependency the extraction
removes; ``$\dyn$/blame'' marks examples that also exhibit the value-as-type
boundary.}
\label{tab:examples}
\Description{A table of eleven example files and the dependent programming idiom each demonstrates.}
\begin{tabular}{@{}llc@{}}
\toprule
File & What it demonstrates & $\dyn$/blame \\
\midrule
\texttt{fin.v}      & bounds-checked indexing; length index erased & \checkmark \\
\texttt{equality.v} & equality indices erased from types; proof terms remain & \\
\texttt{sigma.v}    & dependent pairs; existential over a size & \\
\texttt{printf.v}   & type-safe \texttt{printf}; dependent arity & \\
\texttt{matrix.v}   & dimension-indexed matrices; shape safety & \checkmark \\
\texttt{units.v}    & units of measure; dimension index erased & \checkmark \\
\texttt{avl.v}      & height-indexed balanced trees & \\
\texttt{session.v}  & session-typed protocols; typestate & \\
\texttt{stlc.v}     & intrinsically typed interpreter & \\
\texttt{universe.v} & Tarski universe; a hand-built dynamic type & \checkmark \\
\texttt{functor.v}  & higher-kinded: functors, tagless-final, monads & \\
\bottomrule
\end{tabular}
\end{table}

We highlight three that best illustrate the design.

\paragraph{An intrinsically typed interpreter (\texttt{stlc.v}).}
Object types \texttt{Ty}, contexts \texttt{Ctx}, and well-typed variables
\texttt{Var G t} index the term family \texttt{Tm G t}, so only well-typed object
programs are representable and the evaluator is total by construction. The
identity term and the interpreter extract with all typing scaffolding removed:
\begin{lstlisting}
Extract fun (t : Ty) => tlam emp t t (tvar (snoc emp t) t (vz emp t)).
(* ==>  \t:Ty. (tlam emp t t (tvar (snoc emp t) t (vz emp t))) *)

Extract fun (G : Ctx) (t : Ty) (e : Tm G t) (env : Env G) => eval G t e env.
(* ==>  \G:Ctx. \t:Ty. \e:Tm. \env:Env. (eval G t e env) *)
\end{lstlisting}
\texttt{Tm G t}, \texttt{Env G}, and \texttt{El t} lose their indices; the
interpreter becomes an ordinary evaluator, certified (Theorem~\ref{thm:sim}) to
simulate the dependently typed original.

\paragraph{Units of measure (\texttt{units.v}).}
A quantity \texttt{Qty d} carries a dimension \texttt{d} drawn from an abelian
group; multiplication multiplies dimensions and addition demands equal ones, so
``metres $+$ seconds'' is a type error. Speed extracts with the dimension gone:
\begin{lstlisting}
Extract fun (len time : Dim) (d : Qty len) (t : Qty time) =>
  qmul len (dinv time) d (qinv time t).
(* ==>  \len:Dim. \time:Dim. \d:Qty. \t:Qty. (qmul len (dinv time) d (qinv time t)) *)
\end{lstlisting}
And the family \texttt{Qty} itself, viewed as data (\texttt{fun d => Qty d}),
hits the boundary---emitting $\dyn$ and an internal blame---because it maps a
value to a type.

\paragraph{A hand-built dynamic type (\texttt{universe.v}).}
This example meets the target head on. A universe \texttt{U} of type codes with a
decoder \texttt{El : U -> Set} lets one build \texttt{Dyn}, a dependent pair of a
code and a value it decodes to---literally a dynamically typed value expressed
inside CoC. Codes are ordinary data and extract cleanly; the decoder used as a
function to types is the irreducible residue, extracting to internal blame
through $\dyn$ (\S\ref{sec:overview}). It pinpoints where source and target part
ways and connects the source-level notion of a ``dynamic value'' to the target's
$\dyn$.

%% ==========================================================================
\section{Discussion and Limitations}\label{sec:limitations}
%% ==========================================================================

We are deliberately precise about what is and is not established.

\paragraph{What this paper does not claim.}
This development is not a replacement for Rocq's production extractor, nor is it
a proof-erasing compiler.  The current translation treats proof terms as ordinary
small terms; what it erases are many dependent indices from target types.  The
target typing judgment used by the extraction theorem is not yet a full
kind-regular subject-reduction theorem for the blame calculus.  The simulation
relation used for extraction correctness is syntactic: turning it into a typed
contextual approximation theorem remains open.  These restrictions are
deliberate.  They isolate the central question studied here: how much of a CoC
program's polymorphic skeleton can be preserved in an F$\omega$-style target, and
where must the remaining dependency become dynamic?

\paragraph{Internal blame is reachable by design.}
Theorem~\ref{thm:blamefree} excludes external labels ($\ge 2$) only. The
internal label $\lint$ (identifier $0$) is reachable: a program that exercises a
value-as-type at run time (\S\ref{sec:overview}, \texttt{universe.v}) reduces to
$\blame \lint$. This is intentional---it is the honest image of an untranslatable
construct---and it is exactly why simulation uses $\mathsf{sim\_blame}$: internal
blame approximates any behavior, so it is sound as a placeholder.

\paragraph{The extraction is metatheoretic, not a compiler.}
$\ext$ invokes $\nf$ (normalization) and \texttt{is\_large\_dec} (type-checking)
to make conversion-invariant decisions. It is a total function on derivations
and everything we prove about it holds, but it does not reduce under
\texttt{reflexivity} and is not intended as an efficient back end. Making a
computable, conversion-invariant extraction---perhaps by carrying classification
evidence in a bidirectional elaboration rather than recomputing it---is future
work.

\paragraph{Simulation is syntactic; the bridge to approximation is open.}
Our correctness statements (Theorems~\ref{thm:sim}, \ref{thm:jack}) are in terms
of the syntactic simulation $\sepsim$. We define a contextual approximation
(\texttt{approximates}) in the development, but the containment
$s \sepsim t \Rightarrow s$ approximates $t$ is \emph{not} proved, and for the
current untyped $\sepsim$ it is in fact false: the rule $\mathsf{sim\_left\_tabs}$
relates $\Lambda K.\,s$ to $s$, yet the context $[\cdot]\,[\dyn]$ distinguishes
them (one type-$\beta$-reduces, the other is stuck). A sound bridge would need a
\emph{typed} restriction of $\sepsim$ together with a CIU- or
logical-relations argument. We have since mechanized a target progress
theorem (\S\ref{sec:limitations}), one of the building blocks such an
argument needs, together with the operational determinism already discussed
below; context compatibility of $\sepsim$ under a typed restriction remains
open. We regard $\sepsim$ as a faithful book-keeping device for how source
reduction is mirrored, and closing the gap to observational
approximation---which also requires the target-soundness work
below---as the main open problem.

\paragraph{Operational determinism (mechanized).}
The small-step relation is deterministic: we prove
$\mathsf{step\_deterministic}$, that $[e \reduces e_1]$ and $[e \reduces e_2]$
imply $e_1 = e_2$. The proof rests on three facts we also mechanize: values,
variables, and $\blame$ are irreducible; and distinct ground types are
incompatible ($\mathsf{ground\_compat\_eq}$), which together with a witness
lemma ($\mathsf{compat\_ground\_r\_forces\_arrow}$) makes the ground-cast rules
non-overlapping. This flushed out one real overlap---a cast from $\dyn$ to a
universal type could fire both $\mathsf{step\_generalize}$ and
$\mathsf{step\_collapse}$/$\mathsf{step\_conflict}$---which we fixed by guarding
$\mathsf{step\_generalize}$ with $A \neq \dyn$ (dynamic values reveal their
ground tag first).

\paragraph{Type-level equality and confluence (mechanized).}
The target has genuine F$\omega$ type-level computation: a type-level reduction
$\mathsf{ty\_step}$ with the $\beta$ rule $(\Lambda\alpha{:}K.\,A)\,B \reduces
A[B]$ and congruences, its definitional-equality closure $\mathsf{ty\_equiv}$,
and a typing conversion rule $\mathsf{typing\_conv}$ (with a $\mathsf{wf\_typ}$
premise). We mechanize \textbf{confluence} of $\mathsf{ty\_step}$ by the
parallel-reduction (Takahashi) method---the full de~Bruijn substitution-lemma
chain for $\mathsf{typ}$, a complete-development function, the triangle lemma,
the diamond property, and \textbf{Church--Rosser} for $\mathsf{ty\_equiv}$---and
derive the head-constructor inversions ($\mathsf{ty\_equiv}$ is injective on
$\to$ and $\forall$ and keeps $\dyn$/$\to$/$\forall$/variables distinct). We
also make the $\dyn$-instantiation rules kind-correct: since $\dyn$ lives only at
kind $\kstar$, the reduction/subtyping/compatibility rules that substitute
$\dyn$ into a quantifier are restricted to $\forall\alpha{:}\kstar$.

\paragraph{Target structural metatheory (mechanized).}
On top of confluence we mechanize the structural typing lemmas: \textbf{term
weakening} and \textbf{type weakening} (with the two-namespace context-shift
bookkeeping and $\mathsf{tlift}$-preservation of $\mathsf{compat}$ and
$\mathsf{ty\_equiv}$), \textbf{term substitution} ($\mathsf{typing\_subst}$)
and \textbf{neutral type substitution} ($\mathsf{typing\_tsubst}$, which
requires the substituted type to be $\mathsf{neutral}$, including
$\mathsf{wf\_typ\_tsubst}$ and $\mathsf{compat\_tsubst}$),
\textbf{inversion through} $\mathsf{typing\_conv}$
(recovering a derivation's shape up to conversion for each term former), and
\textbf{canonical forms} (a value at an arrow/$\forall$ type is a
$\lambda$/$\Lambda$). All of these are admit-free.

\paragraph{Target progress (mechanized).}
Building on the structural metatheory above and on executable
$\mathsf{compat}$, we prove $\mathsf{progress}$: every closed term well typed
under the permissive judgment is a value, steps, or is $\blame p$ for some $p$.
The cast case is exactly the case analysis on $\mathsf{compat}(A,B)$ sketched
above: every constructor of the executable judgment corresponds to a concrete
reduction rule (ID, WRAP, the structural $\forall/\forall$ step
$\mathsf{step\_all\_all}$, GENERALIZE, INSTANTIATE, GROUND, COLLAPSE/CONFLICT),
so there is no leftover ``stuck cast'' shape to handle by fiat---progress falls
directly out of $\mathsf{compat\_dec}$. The two remaining new lemmas are:
head-distinctness for $\mathsf{ty\_equiv}$ extended to $\mathsf{tvar}$ and
$\Lambda$ heads (ruling out a value at a variable or type-abstraction type by
canonical forms), and $\mathsf{is\_gnd}$/$\nu$-tamper detection generalized from
the bare type-variable case to \emph{any} neutral ground tag $N$
(\S\ref{sec:target}): a sealed variable can now hide inside $N\,A$, not just a
bare $\alpha$, so the tamper guards inspect $\mathsf{neutral}$ rather than
pattern-matching on $\mathsf{tvar}$.

\paragraph{Kind-regular target typing is future work.}
The typing judgment used by Theorem~\ref{thm:welltyped-extraction} is permissive:
it does not embed $\mathsf{wf\_typ}$ (and $\mathsf{wf\_ground}$) premises at every
constructor, so it certifies term formation rather than full kind-regularity. The
\emph{extraction side} of the eventual kind-regular judgment is, however, already
mechanized: both premises the large-\textsc{App} conversion needs are proved---
kind regularity of type extraction (\texttt{extract\_typ\_wf\_sort}, that
$\extt(\Gamma,T)$ is well kinded at $\kstar$) and the conversion lemma
(\texttt{extract\_typ\_tsubst\_coc\_equiv}). Getting kind regularity right was
subtle: the naive kind index ($\extt(\Gamma,T)$ has kind $\extk(U)$ for $T$'s CoC
type $U$) is \emph{false} (for an application whose codomain mentions its bound
variable, $\extk$ does not commute with the source substitution), so the proof
threads the target kind through the CoC \emph{kinding} derivation, restricted to
genuinely type-level subterms. What remains is to \emph{rewire} the target term
typing judgment to embed these kinding premises (and a kind-aware
$\mathsf{wf\_ground}$ over neutral ground tags).

\paragraph{Target preservation is future work.}
We prove progress for the current judgment (Theorem via \texttt{progress}), but a
target subject-reduction theorem should be proved only \emph{after} the
kind-regular judgment above is introduced---preservation of a non-kind-regular
judgment is not the meaningful final theorem. The structural metatheory it will
build on is already in place and admit-free (term/type weakening, term
substitution, neutral type substitution, inversion through conversion, canonical
forms, and the $\mathsf{ty\_equiv}$/$\mathsf{typing\_conv}$ equational theory).
None of this affects the current blame-safety results, which are stated for the
raw syntactic safety predicates rather than for a kind-regular target typing
theorem.

\paragraph{$\dyn$-freedom is fragment-restricted.}
Theorem~\ref{thm:dynfree} gives an exact iff characterization for the
source-type skeleton fragment $\mathsf{fo}$. Full CoC types outside this
fragment can extract to types containing $\dyn$; whether a larger $\dyn$-free
fragment exists (e.g.\ admitting higher kinds) is open.

\paragraph{Inductive families are axiomatized.}
The bare PTS has no inductive types, so our examples model families and their
eliminators as \texttt{Axiom}s. This is standard for a PTS presentation and does
not affect the metatheory (axioms are just context bindings), but extending the
source to the Calculus of Inductive Constructions---with native inductives and
their large eliminations---would let programs that \emph{compute} types from data
appear, stressing the value-as-type boundary in richer ways.

%% ==========================================================================
\section{Related Work}\label{sec:related}
%% ==========================================================================

\paragraph{From CoC proofs to F$\omega$ programs.}
The closest prior work is Paulin-Mohring's extraction of
F$\omega$ programs from proofs in the Calculus of Constructions
\cite{paulinmohring1989fomega}.  Paulin-Mohring uses realizability and a
distinction between informative and non-informative propositions to remove the
logical part of a development; the extracted programs are non-dependent
polymorphic terms.  We build on the same insight that a substantial CoC
fragment has an F$\omega$ skeleton.  Our contribution is different: we give a
Rocq mechanization whose target is an explicit blame calculus, not pure
F$\omega$, and we use $\dyn$ and internal blame to make the non-F$\omega$
residue explicit.  Thus our extraction is optimistic and total for the
formalized CoC core, while Paulin-Mohring's extraction is a realizability-based
program extraction with proof/logical erasure.

\paragraph{Rocq extraction and verified extraction.}
Rocq's production extraction emits ML-family code---currently OCaml, Haskell, or
Scheme---and erases dependent/proof structure that those target languages cannot
type directly \cite{letouzey2003extraction,letouzey2004these}; erasure is guided
by sort information much as our large/small split is.  Forster, Sozeau, and
Tabareau verify a modern extraction pipeline from Coq to
OCaml~\cite{forster2024verified}.  Because OCaml lacks a formal semantics, their
verified core targets Malfunction, a formal specification of the OCaml
compiler's Lambda intermediate language.  The released plugin produces
\emph{untyped} Malfunction programs together with typed interface files, and the
main guarantees are operational correctness and safe interoperability for
first-order data, not well-typedness of a residual target term.  Our theorem is
therefore orthogonal: we do not target production OCaml, but we do prove that
the extracted residual term is well typed in an explicit polymorphic blame
calculus.  The two results are strong in different dimensions: theirs gives
operational correctness and safe interoperability for a production language;
ours gives well-typedness in a rich polymorphic target at the cost of
practicality.

\paragraph{Verified erasure and verified compilation.}
MetaCoq/MetaRocq verifies type checking and type/proof erasure for Coq's kernel,
including erasure to an untyped lambda calculus
\cite{sozeau2020metacoq,sozeau2020coqcoqcorrect}.
CertiCoq \cite{anand2017certicoq} verifies a compiler from Gallina toward
lower-level executable code, and \OE{}uf \cite{mullen2018oeuf} minimizes the
extraction TCB through an untyped intermediate language.  Erasure in dependently
typed \emph{programming} languages, such as Idris'
quantitative/erasure analysis \cite{brady2013idris,tejiscak2020erasure} and
erasure in pure type systems \cite{mishra2008erasure}, removes computationally
irrelevant arguments while staying within a dependent target.  These systems aim
to produce executable code with small trusted bases.  We instead isolate a typed
intermediate-language question: after dependent information is removed from
types, what polymorphic structure can still be expressed statically, and how
should the remaining dependency be represented dynamically?

\paragraph{Gradual typing, blame, and polymorphism.}
Gradual typing \cite{siek2006gradual,siek2015refined} interleaves static and
dynamic checking; blame \cite{findler2002contracts,wadler2009blame} attributes
run-time cast failures to boundaries. Ahmed et al.'s \emph{Blame for All}
\cite{ahmed2011blame} adds polymorphism, with the Jack-of-All-Trades principle
and dynamic sealing. We adapt its $\nu$-style sealing and
blame-safety structure to an F$\omega$-style target.  Two caveats are important.
First, the original Blame-for-All artifact reports that the published proof of
the Jack-of-All-Trades Principle is wrong, so we treat our Jack theorem only as
the mechanized syntactic simulation lemma that the code actually proves.  Second,
where Blame for All's ground predicate for System~F is restricted to a bare
type-variable tag, we generalize to \emph{neutral} F$\omega$ proper types
($N ::= \alpha \mid N\,A$, \S\ref{sec:target}), so a type-family application at
the boundary is also directly ground; kind-regularity of the target typing
judgment remains future work (\S\ref{sec:limitations}).
Abstracting Gradual Typing \cite{garcia2016abstracting} and gradual type theory
\cite{new2019gradual,new2018gradualdt} give systematic accounts of gradual
languages and their metatheory.

\paragraph{Gradual dependent types.}
A line of work gradualizes dependent types directly: approximate normalization
\cite{eremondi2019approximate} and Gradual CIC \cite{lennon2022gradual} equip a
dependent theory with a dynamic type internally, and Tanter and Tabareau bring
gradual certified programming to Coq \cite{tanter2015gradualcertified}. These
make the \emph{source} gradual. We instead keep the source fully static (CoC) and
make the \emph{target} gradual, using $\dyn$ solely to absorb the dependency that
the polymorphic target cannot represent---a translation, not a language design.
Cayenne \cite{augustsson1998cayenne} compiled a dependently typed language by
erasing to an untyped core; intrinsically typed interpreters (our
\texttt{stlc.v}) are folklore in Agda \cite{norell2008agda} and Idris.  Our
contribution is not the source programs but the verified image of such programs
in a typed gradual calculus.

%% ==========================================================================
\section{Conclusion}\label{sec:conclusion}
%% ==========================================================================

We have shown that extraction from a dependently typed language need not throw
away all types: with a gradual target, a single verified translation preserves
the parametric-polymorphic skeleton of CoC programs, removes many dependent
indices from target types, and confines target-inexpressible type-level residue---most visibly values used as
types---to $\dyn$ and, when forced operationally, to reserved internal blame.  In this sense the work returns to Paulin-Mohring's
CoC-to-F$\omega$ vision~\cite{paulinmohring1989fomega}, but with a gradual
target and a mechanized blame-safety story.  The whole pipeline, including a
machine-checked F$\omega$-style polymorphic blame calculus, is formalized in
Rocq, with the target and all of the extraction's metatheorems---well-typed
extraction, simulation, and derivation independence---fully admit-free
(\S\ref{sec:mech}); a REPL with verified name
recovery makes its behavior
inspectable on realistic examples. The current extractor does not implement Coq-style proof
erasure; proof terms remain ordinary residual terms. The results chart a design
point between fully typed and fully untyped extraction; the natural next steps
are closing the extraction weakening/substitution commutation gap, a
kind-regular target type-safety theorem, the bridge from simulation to
observational approximation, a proof-erasing optimization pass, and an extension
to inductive families and their large eliminations.

\bibliographystyle{ACM-Reference-Format}
\bibliography{references}

\end{document}
